\section{Pemrosesan Form dan OOP Dasar}

PHP memproses form dengan membaca \texttt{\$\_POST} atau \texttt{\$\_GET}, memvalidasi input, dan memproses data (simpan ke basis data, kirim email, dll.) \cite{php-manual}. Selalu gunakan \texttt{htmlspecialchars()} atau escape saat menampilkan data pengguna untuk mencegah XSS \cite{php-right-way}.

PHP mendukung pemrograman berorientasi objek (OOP): class, object, property, method, constructor, visibility (\texttt{public}, \texttt{protected}, \texttt{private}) \cite{php-manual}. OOP membantu mengorganisasi kode aplikasi yang besar; misalnya class untuk koneksi database atau model data \cite{php-right-way}.

\begin{lstlisting}[caption={Contoh PHP memproses form}, basicstyle=\ttfamily\small, frame=single]
<?php
if ($_SERVER['REQUEST_METHOD'] === 'POST') {
  $nama = htmlspecialchars(trim($_POST['nama'] ?? ''));
  if ($nama !== '') echo "Halo, " . $nama;
}
?>
\end{lstlisting}
